\documentclass[pdflatex,sn-mathphys-num]{sn-jnl}

\usepackage{graphicx}%
\usepackage{multirow}%
\usepackage{amsmath,amssymb,amsfonts}%
\usepackage{amsthm}%
\usepackage{mathrsfs}%
\usepackage[title]{appendix}%
\usepackage{xcolor}%
\usepackage{textcomp}%
\usepackage{manyfoot}%
\usepackage{booktabs}%
\usepackage{algorithm}%
\usepackage{algorithmicx}%
\usepackage{algpseudocode}%
\usepackage{listings}%
\usepackage{lmodern}%
\usepackage{siunitx}
\usepackage{microtype}
\usepackage{derivative}
\usepackage[version=4]{mhchem}
\usepackage{cleveref}
\usepackage[spanish]{babel}

\renewcommand{\arraystretch}{1.2}
\sisetup{group-digits=true,
  group-separator={\,},
  separate-uncertainty
}

\theoremstyle{thmstyleone}%
\newtheorem{theorem}{Teorema}%
\newtheorem{proposition}[theorem]{Proposición}%

\theoremstyle{thmstyletwo}%
\newtheorem{example}{Ejemplo}%
\newtheorem{remark}{Observación}%

\theoremstyle{thmstylethree}%
\newtheorem{definition}{Definición}%

\DeclareSIUnit\angstrom{\text{\r{A}}}

\raggedbottom
\unnumbered% uncomment this for unnumbered level heads

\begin{document}

\title[Machine Learning Científico]{Aprendizaje Automático Científico y
Dinámica No Lineal: Resumen de la Charla de Carlos Andrés Gómez Vasco}

\author{\fnm{Julian L.}
\sur{Avila-Martinez}}\email{jlavilam@udistrital.edu.co}

\affil{\orgdiv{Programa Académico de Física}, \orgname{Universidad Distrital
Francisco José de Caldas}, \orgaddress{\city{Bogotá D.C.},
\country{Colombia}}}

\abstract{
  Se expone el cambio de paradigma introducido por el aprendizaje automático
  científico como puente entre la simulación física estricta y la inferencia
  empírica de datos. Se analiza la incorporación de simetrías y leyes de
  conservación en arquitecturas computacionales mediante funcionales de costo
  modificados, mitigando la dependencia absoluta de volúmenes masivos de datos.
  Adicionalmente, se examina el aprendizaje de operadores para mapear dinámicas
  caóticas en espacios vectoriales de comportamiento lineal y su aplicación en
  la construcción de gemelos digitales para la predicción en tiempo real.
}

\keywords{Aprendizaje automático científico, redes neuronales, regresión
simbólica, isomorfismos no lineales, gemelos digitales}

\maketitle

\section{Epistemología Computacional y Redes Informadas}
\label{sec:epistemologia}

La formulación de métodos numéricos tradicionales exige mallas espaciales
densas para resolver sistemas de ecuaciones diferenciales, garantizando un
seguimiento riguroso de cada variable bajo un alto costo computacional. En
riguroso contraste, el aprendizaje automático empírico opera como un
mecanismo opaco desprovisto de restricciones físicas, requiriendo volúmenes
masivos de datos limpios para establecer correlaciones. El aprendizaje
automático científico, o SciML, resuelve esta divergencia embebiendo
aserciones analíticas directamente en la arquitectura del algoritmo. Esta
restricción física impuesta disminuye drásticamente el espacio de fases
accesible, permitiendo la convergencia estricta del modelo aun con conjuntos
de datos limitados o con ruido estocástico.

La materialización de este esquema se evidencia en las redes neuronales
informadas por la física, conocidas en la literatura como PINNs. Estas
estructuras imponen penalizaciones analíticas en el funcional de costo
mediante el uso sistemático de la diferenciación automática, castigando
matemáticamente aquellas configuraciones espaciales que violen principios
básicos como la conservación de la masa o el momento. Aunque estas redes
prescinden de la discretización topológica mediante mallas, exhiben un
costo de entrenamiento sustancial y presentan dificultades de convergencia
ante discontinuidades severas, fenómenos característicos de la dinámica de
fluidos.

\section{Descubrimiento Simbólico y Endomorfismos}
\label{sec:descubrimiento}

La inferencia de las leyes rectoras subyacentes a partir de mediciones se
formaliza a través de la regresión simbólica impulsada por algoritmos
genéticos. El método inicializa un diccionario exhaustivo de funciones y
operadores abstractos, aplicando transformaciones algebraicas sucesivas a
las expresiones postuladas. El algoritmo converge al identificar la estructura
algebraica más simple que determina invariablemente la evolución temporal del
sistema, operando en esencia como un mecanismo de ingeniería inversa sobre
las leyes empíricas.

En el régimen de la dinámica caótica y fuertemente no lineal, la frontera
teórica del SciML abandona la búsqueda de soluciones particulares para
enfocarse en el aprendizaje de representaciones algebraicas del operador. Las
redes profundas mapean las funciones de estado del sistema hacia un nuevo
espacio vectorial abstracto. Bajo este difeomorfismo, la evolución del
sistema asume un comportamiento estrictamente lineal. Este isomorfismo elude
las limitaciones intrínsecas del cálculo clásico sobre variedades complejas
y garantiza que la representación del endomorfismo se generalice para
resolver múltiples clases de sistemas físicos isomorfos.

\section{Dinámica Predictiva y Gemelos Digitales}
\label{sec:gemelos}

La síntesis de arquitecturas robustas de SciML viabiliza la ejecución de
gemelos digitales, conceptualizados como proyecciones virtuales e
interactivas de los sistemas físicos. Al suministrar flujos de datos
sensoriales incompletos, tales como las fluctuaciones térmicas en celdas
de baterías de vehículos eléctricos, el modelo predice gradientes y estados
coherentes con las restricciones termodinámicas del medio. La aceleración
computacional que proporciona la formulación sobre el cálculo convencional
permite optimizar la funcionalidad estructural en tiempo real, previniendo
disipaciones o fallos catastróficos del sistema.


\end{document}
